Source : OCDE, Revenue Statistics (jeu DSD\_REV\_COMP\_OECD@DF\_RSOECD, rubrique 1200, impôts sur le revenu, les bénéfices et les gains en capital des sociétés, toutes administrations publiques), téléchargé le 16/11/2026 ; extraction : fiscalite/is/donnees/extraire\_ocde\_is\_pib.py.\par Note : année 2023, la dernière disponible pour les 38 pays. Médiane de l'OCDE : 3,41 \%. Pour la France, 2,85 \% en 2022 (année de bénéfices exceptionnels) et 2,25 \% en 2024 (31e sur 36). La contribution exceptionnelle sur les grands groupes, perçue en 2025 et 2026, n'est pas encore dans ces chiffres. Plusieurs pays du haut du classement doivent leur rang à des ressources naturelles ou à des bénéfices de multinationales localisés chez eux.
